\section{Reproducibility}
\label{sec:reproducibility}

The code, the committed documents of case C16 and every stored result are in the repository named in the header of the first page. All commands run from the repository root with \texttt{python} on the environment of the product (the \texttt{riskvalidation} package 1.1.0 with the market, machine-learning and stress extras, and the requirements of the case pipeline). The experiments read the committed documents of \texttt{data/derived/C16}; no raw data are needed, and no path of a machine is written in any file.

\begin{itemize}
\item \texttt{experiments/run\_reverse\_c16.py} rebuilds the three bases, calls the engine, and writes \texttt{results\_reverse\_c16.json} (reproductions, derivative, convention band, labelled run).
\item \texttt{experiments/run\_corpus.py} generates the five corpora from seeds 20261107 to 20261111, runs the oracle, the solver and the comparators, and writes \texttt{results\_corpus.json}; it took \kWallMin{} minutes on one machine with other jobs running (\kProc{}, Python \kPy{}).
\item \texttt{experiments/run\_theory.py} writes \texttt{results\_theory.json} from seed 20261120.
\item \texttt{experiments/hypotheses.json} holds the predictions and the dates at which they were written; \texttt{experiments/test\_core.py} and \texttt{experiments/test\_results.py} test the shared code and the stored results.
\item \texttt{make\_numbers.py}, \texttt{make\_tables.py} and \texttt{make\_figures.py} write \texttt{tables/numbers.tex}, the tables and the vector figures from the stored results, so that no number of the text is typed; \texttt{fetch\_references.py} writes \texttt{references.bib} from the records of the cited works.
\end{itemize}

The times in \cref{tab:main-results} are medians of measured calls, written by the experiment scripts to the stored results; they depend on the machine and are not a property of the method. The solve on the Adverse base takes \timeAdvMs{} ms and on the Extreme base \timeExtMs{} ms.
